\originalchapter{粒子、场与对称性}\label{ch:symmetry}
\originalsection{粒子分类与实验身份}
基本费米子包括夸克与轻子, 自旋为 $1/2$. 夸克带颜色并组成强子: 介子常可用夸克--反夸克描述, 重子常用三个夸克描述; 更一般复合强子不被这两个最简单组态穷尽. 光子、胶子、$W^\pm,Z$ 是自旋 $1$ 规范玻色子, Higgs 是自旋零粒子. 质子是复合费米子, pion 是复合玻色子, ``基本''与``自旋统计类型''是两种分类.

粒子可由质量峰、寿命、衰变模式和量子数识别. 窄共振的 Breit--Wigner 分母为 $(s-M^2)^2+M^2\Gamma^2$, $\Gamma$ 反映开放通道的总宽度. 由 $\tau=\hbar/\Gamma$ 将寿命与能量宽度对应. 自由场的产生、湮灭描述让粒子数可改变; 本分编只使用足够计算过程的场与顶点, 完整量子化和重整化证明由独立 QFT 卷承接.

\originalsection{连续对称与 Noether 流}
拉格朗日密度 $\mathcal L(\phi,\partial_\mu\phi)$ 的 Euler--Lagrange 方程为
$\partial_\mu[\partial\mathcal L/\partial(\partial_\mu\phi)]-\partial\mathcal L/\partial\phi=0$. 若无穷小变换 $\delta\phi=\epsilon\Delta\phi$ 只使 $\delta\mathcal L=\epsilon\partial_\mu K^\mu$, 则
\begin{equation}
 j^\mu=\frac{\partial\mathcal L}{\partial(\partial_\mu\phi)}\Delta\phi-K^\mu,\qquad\partial_\mu j^\mu=0.
\end{equation}
多分量场对各独立分量求和. 证明只需把拉格朗日变化写成方程运动项加总散度; 在壳上前项为零. 空间积分 $Q=\int j^0\dd^3x$ 在边界无净流时守恒.

复标量全局相位 $\phi\to\e^{\ii\epsilon}\phi$ 给电荷流. 若把相位推广为任意位置函数, 普通导数会带 $\partial_\mu\epsilon$, 必须引入规范场及协变导数. 全局对称与局域规范冗余概念有别, 不应把每个局域变换都当作一种可独立测量的新守恒荷.

\originalsection{同位旋与离散变换}
忽略 $u,d$ 质量差和电磁作用时, 强相互作用近似具有同位旋对称, $(p,n)$ 是 $I=1/2$ 双重态, pion 构成 $I=1$ 三重态. 第三分量 $I_3$ 由粒子在多重态中的位置指定, 不是空间自旋分量. 电磁作用区别 $u,d$ 电荷, 会破坏完整同位旋; 但它保持各夸克味的净数, 因而通常仍保持相应 $I_3$.

宇称 $P$ 令 $\mathbf r\to-\mathbf r$, 普通极矢量 $\mathbf p,\mathbf E$ 变号, 轴矢量 $\mathbf L,\mathbf S,\mathbf B$ 不变. 时间反演 $T$ 令 $t\to-t$, 速度、动量、自旋和磁场变号, 位置和静电场不变. 因此 $\mathbf S\cdot\mathbf E$ 在 P 与 T 下都奇, 永久电偶极矩与这些对称性破坏有关; $\mathbf S\cdot\mathbf B$ 在两者下都偶.

$C$ 将粒子换成反粒子, CP 联合电荷共轭与宇称. QED 和通常的强相互作用近似保持 P 与 CP, 弱相互作用违反 P, 还可经 CKM 相位违反 CP. 强 CP 是否破坏取决于 $\theta$ 项; ``强作用一定严格 CP 不变''是超出实验与理论条件的绝对说法.

\originalsection{选择规则与近似的层次}
能动量、电荷和角动量守恒是判断过程的第一关. 味量子数如奇异数在强、电磁过程中守恒, 在带电弱流中可以变化. 单个轻子味数因中微子混合不严格守恒; 低能标准模型通常讨论的总轻子数也会受到可能的 Majorana 质量和非微扰效应限制. 课程入门表必须注明层次, 避免把方便的近似说成所有能量和所有扩展理论中的绝对定律.

交换相同费米子外腿时振幅相对负号来自反对易, 不是任选的图形方向约定. 全局整体相位不影响概率, 两个图之间的相对相位则决定干涉. 下一章把这些规则落实到 toy theory.
